\subsection{无杠杆约束的经济含义与桥水框架的区别}
\label{sec:unlevered}
本文选择多头、满额投资且无杠杆的预算集合，是为了研究境内上市ETF能够实现的资产配置。这里的“工具较少”具体指本文机会集合的覆盖范围：在具有连续历史、可统一核算的30只工具中，权益品种较多，债券风险主要由境内名义国债、信用债与可转债承担，通胀相关暴露主要通过商品类工具表达。与全球机构投资者可组合的跨国名义债、通胀挂钩债、货币对冲和衍生品体系相比，这一集合覆盖的独立风险来源及风险调节方式较为有限。因此，本文把资本在现有ETF之间的分配作为完整研究对象。

桥水对All Weather发展过程的官方说明强调经济增长与通胀意外的平衡，并讨论利用期货、全球通胀挂钩债和货币对冲组织风险暴露，以及借助杠杆调整低风险资产的风险规模\cite{bridgewater2012allweather}。本文借鉴风险分散的思想，将其落实到境内ETF的资本预算约束。两者所用工具、融资条件和风险映射不同，本文的Global RRP应按其自身可行集合评价。

无杠杆约束使风险预算与资本占用直接相连。低波动资产若要承担较多风险，通常需要配置较高资本权重；当总权重固定为一时，其余资产可用的资本相应减少。主模型的收益目标因此取当期可行参考组合的预测收益，使收益要求与现有工具和资本约束相协调。相同逻辑也决定了辅助比较的设计：通过货币ETF降低简单组合的风险，而将混合比例限制在零到一之间。

杠杆能够扩大低波动资产的风险暴露，也会引入融资成本与保证金需求。Brunnermeier 与 Pedersen 在《Review of Financial Studies》解释了市场流动性与融资流动性的相互强化：融资条件收紧可能促使持仓收缩，并进一步影响交易条件\cite{brunnermeier2009liquidity}。因此，在本文的ETF价格与固定费率数据框架内，坚持无杠杆能够使收益、风险和费用来自同一组可核算持仓。本文完成的是给定资本预算下的风险收益配置，后续引入融资工具时，应将融资利率、保证金、展期和压力状态下的现金需求同时纳入。
